\section{Calibration du facteur taux}
\label{sec:taux}

\subsection{Objectif et jeux de calibration}

Comme \citet[section~4.2.7]{BrigoMercurio2006} et \citet[éq.~55]{Russo}, on minimise les écarts relatifs entre
les prix du modèle et les prix de marché de $N$ swaptions à la monnaie :
\begin{equation}
\min_{a_x,\sigma_x}\ \sqrt{\sum_{i=1}^{N}\Big(\frac{\mathit{Swpt}_i(a_x,\sigma_x) - \mathit{Swpt}_i^M}{\mathit{Swpt}_i^M}\Big)^2}.
\end{equation}
Les prix de marché sont les prix de Bachelier ; les prix du modèle sont les prix exacts de Jamshidian. Comme le
prix d'une swaption à la monnaie est linéaire en volatilité normale, l'erreur relative de prix est aussi l'erreur
relative de volatilité, et les résidus s'expriment en points de base de volatilité normale. L'optimisation
(L-BFGS-B, $a_x \in [0{,}001\,;\,1]$, $\sigma_x \in [0{,}01\,\%\,;\,5\,\%]$) part de plusieurs points, et l'on
retient le meilleur.

Calibré sur la surface complète (291 swaptions, en excluant celles dont le swap se termine au-delà de cinquante
ans), le modèle donne $a_x = 0{,}0136$ et $\sigma_x = 0{,}735\,\%$, pour une erreur quadratique moyenne de 4,5~pb.
La carte des résidus (figure~\ref{fig:calibration_taux}, à gauche) montre une erreur structurelle : avec deux
paramètres constants, le modèle surévalue les volatilités des expiries courtes, jusqu'à six mois, et sous-évalue
celles des expiries de un à cinq ans sur les tenors courts.

On calibre donc sur les swaptions qui comptent pour le produit, comme le recommandent
\citet[section~4.2.7]{BrigoMercurio2006} et \citet[section~6.3.7]{Chaix2021} : la diagonale co-terminale
$T_0 + n = 20$ ans, vingt ans étant la maturité résiduelle de l'OAT sous-jacente arrondie à l'année. Les swaptions
dont le swap se termine à la maturité de l'obligation portent l'information sur la volatilité des taux qui
compte pour l'option. Le cube ne cote que les tenors 1 à 10, 15, 20, 25 et 30 ans ; les tenors intermédiaires de
la diagonale sont obtenus par interpolation linéaire de la volatilité normale en tenor, à expiry donnée. Avec
$T_0 \ge 1$ an, la diagonale compte douze swaptions d'expiry 1 à 15 ans, dont huit interpolées.

\subsection{Résultats}

\begin{table}[htbp]
\centering\small
\caption{Calibration de $(a_x, \sigma_x)$ sur les diagonales co-terminales (prix exacts de Jamshidian)}
\label{tab:diagonales}
\begin{tabular}{@{}lcccc@{}}
\toprule
Diagonale & Swaptions (dont interpolées) & $a_x$ & $\sigma_x$ (\%) & Erreur relative (\%) \\
\midrule
10 ans & 9 (0) & 0,0010 & 0,7235 & 1,39 \\
\textbf{20 ans} & \textbf{12 (8)} & \textbf{0,0107} & \textbf{0,7367} & \textbf{1,44} \\
30 ans & 14 (9) & 0,0321 & 0,9453 & 1,62 \\
20 ans, expiries 1 à 10 ans & 10 (8) & 0,0010 & 0,6761 & 1,27 \\
\bottomrule
\end{tabular}
\end{table}

Sur la diagonale 20 ans, $(a_x,\sigma_x)$ sont identifiés conjointement : l'objectif n'a qu'un minimum, que les
départs proches atteignent. On retient
\begin{equation}
a_x = 0{,}01067,\qquad \sigma_x = 0{,}7367\,\%.
\end{equation}
Le plus grand résidu, 2,3~pb, porte sur l'expiry d'un an, dont la volatilité de marché est interpolée ; les
autres ne dépassent pas 1,3~pb (figure~\ref{fig:calibration_taux}, à droite). Les expiries de 12 et 15 ans
dépassent la plus longue échéance d'option, mais ce sont elles qui identifient la vitesse de retour : restreinte
aux expiries de 1 à 10 ans, la diagonale ne contraint plus $a_x$, et l'objectif décroît jusqu'à la borne
inférieure (tableau~\ref{tab:diagonales}). Les diagonales 10 et 30 ans donnent d'autres paramètres : avec deux
paramètres constants, les paramètres de taux dépendent du sous-jacent, ce qui justifie de calibrer sur la
diagonale de l'OAT évaluée.

\begin{figure}[htbp]
\centering
\includegraphics[width=\textwidth]{calibration_taux}
\caption{Calibration du facteur taux. À gauche, résidus du modèle calibré sur la surface complète (modèle moins
marché, en points de base de volatilité normale) ; à droite, diagonale co-terminale 20 ans, marché et modèle
calibré.}
\label{fig:calibration_taux}
\end{figure}

La vitesse de retour calibrée, 0,011, est très inférieure à l'ordre de grandeur usuel de 0,2 à 2
\citep[CM2, p.~6]{Hillairet2026}. C'est la diagonale qui l'impose ; la section~\ref{sec:resultats} montre que le
prix de l'option en dépend peu tant que cette diagonale reste reproduite.

Enfin, la diagonale permet de mesurer l'approximation à poids gelés en un facteur, où le prix exact est connu.
L'écart entre la formule de Black à poids gelés et Jamshidian ne dépasse pas 0,07~\% du prix, et une calibration
sur les prix à poids gelés donnerait $a_x = 0{,}01051$ et $\sigma_x = 0{,}7352\,\%$, des valeurs presque
identiques.
